\documentclass[aps,pra,reprint,groupedaddress,twocolumn,superscriptaddress,linenumbers]{revtex4-2}

\usepackage{graphicx}% Include figure files
\usepackage{dcolumn}% Align table columns on decimal point
\usepackage{bm}% bold math
\usepackage{amsmath}
\usepackage{siunitx}
\usepackage{enumitem}
\usepackage{booktabs}

\begin{document}

\preprint{APS/123-QED}

\title{Quantum-Traceable Refractometry at 1762 nm for Barium-Ion Quantum Networks}

\author{Wei Wu}
\affiliation{Albert-Ludwigs-Universität Freiburg, Physikalisches Institut, Hermann-Herder-Straße 3, 79104 Freiburg, Germany}
\affiliation{EUCOR Centre for Quantum Science and Quantum Computing, University of Freiburg, Hermann-Herder-Straße 3, 79104 Freiburg, Germany}
\author{Tobias Schätz}
\affiliation{Albert-Ludwigs-Universität Freiburg, Physikalisches Institut, Hermann-Herder-Straße 3, 79104 Freiburg, Germany}
\affiliation{EUCOR Centre for Quantum Science and Quantum Computing, University of Freiburg, Hermann-Herder-Straße 3, 79104 Freiburg, Germany}
\author{Ulrich Warring}
\affiliation{Albert-Ludwigs-Universität Freiburg, Physikalisches Institut, Hermann-Herder-Straße 3, 79104 Freiburg, Germany}

\date{\today}

\begin{abstract}
Precise knowledge of the atmospheric refractive index at the \SI{1762}{\nano\meter} barium-ion transition wavelength is critical for environmental compensation in quantum networks. This work establishes a quantum-traceable measurement framework that directly links ambient refractive index to a single trapped $^{138}$Ba$^+$ ion frequency standard. Using a dual-wavelength interferometer and comprehensive statistical validation, we determine refractive index coefficients with part-per-billion precision: $\alpha_T = \SI{-8.8474e-7}{\per\degreeCelsius}$, $\alpha_H = \SI{-1.3152e-8}{\per\percent}$, and $\alpha_P = \SI{+2.5949e-7}{\per\hecto\pascal}$. The measurement reveals a \SI{12.4}{\percent} enhanced humidity sensitivity compared to standard atmospheric models, consistent with Kramers-Kronig predictions near water absorption lines. A complete uncertainty budget identifies environmental sensor calibration as the dominant error source ($\delta n \sim 2.12\times10^{-7}$). These results provide the metrological foundation for quantum network operation under realistic environmental conditions.
\end{abstract}

% \pacs{42.50.-p, 42.25.Bs, 07.60.Ly, 03.67.Hk}

\maketitle


\section{Introduction}

A pivotal challenge in moving quantum technologies from controlled laboratories to real-world applications is their sensitivity to environmental fluctuations. For quantum networks based on trapped ions like $^{138}$Ba$^+$, phase stability over optical interconnects at specific wavelengths, such as 1762 nm, is critical. This stability is compromised by atmospheric refractive index changes, which are typically compensated for using generalized models like the Ciddor equation. These models originate from atmospheric science and classical metrology, where they are validated for broad spectral ranges and climatological averages. Concurrently, the field of quantum information and control has achieved remarkable isolation and coherence for individual qubits, while quantum metrology has established frequency standards with fractional uncertainties at the $10^{-18}$ level. However, these ultraprecise quantum capabilities have not been translated to the task of measuring the dynamic environmental parameters that now limit system performance.

This divergence in capability has created a "traceability gap." Quantum systems are engineered toward fundamental noise limits, yet they operate within an environment whose characterization lacks direct traceability to a quantum standard. The result is that key parameters, such as the refractive index at a qubit's specific optical wavelength, carry an uncertainty that is not commensurate with the qubit's intrinsic precision. For a $^{138}$Ba$^+$ ion network node, this gap means the coefficients required for precise phase compensation at 1762 nm cannot be derived from the ion's own frequency reference. Instead, reliance on indirect models introduces a dominant, unquantified source of error that constrains entanglement fidelity and network scalability.

In this Letter, we bridge this gap by creating a sensor that provides quantum traceability for an environmental variable. We demonstrate a refractometer for 1762 nm whose measurement is directly traceable to the transition frequency of a single trapped $^{138}$Ba$^+$ ion. This is achieved through a co-designed apparatus that uses the ion as a primary standard to calibrate a dual-wavelength interferometer, thereby establishing a continuous SI-traceable chain to the ambient refractive index. Our measurement yields the modified Edlén coefficients for 1762 nm with an uncertainty budget anchored to the quantum standard. This provides the quantum information and control community with a deterministic tool for environmental compensation, directly addressing a key performance bottleneck. For metrology, this work demonstrates a pathway to extend quantum standards beyond traditional quantities like time and frequency, into the domain of environmental sensing. For atmospheric and environmental science, it provides a method for validating fundamental refractive index models at a specific, technologically critical wavelength with unprecedented precision. By integrating these perspectives, our approach enables a new paradigm of quantum-aware environmental characterization, which is essential for building scalable and robust quantum technologies.

\section{System Architecture}

The measurement system integrates quantum frequency references with a modified NIST LM10 Michelson interferometer \cite{fox1999reliable}. A \SI{780}{\nano\meter} reference laser is stabilized via saturation absorption spectroscopy of $^{87}$Rb, providing $\delta f \sim \SI{100}{\kilo\hertz}$ stability. The \SI{1762}{\nano\meter} probe laser is stabilized to an ultra-low expansion (ULE) cavity that is referenced to a single trapped $^{138}$Ba$^+$ ion, achieving $\delta f < \SI{1}{\kilo\hertz}$. This quantum reference contributes $\delta n < 10^{-12}$ uncertainty to the refractive index measurement.

The interferometer operates with \SI{4.2}{\meter} optical path difference and air-floated mirror stages for vibration isolation. Simultaneous fringe counting at both wavelengths enables common-mode rejection of environmental noise. Co-located ambient sensors (BME280) measure temperature, humidity, and pressure with specified uncertainties of \SI{\pm 1}{\degreeCelsius}, \SI{\pm 3}{\percent} RH, and \SI{\pm 0.6}{\hecto\pascal}, respectively. These environmental sensors represent the dominant error source in the measurement chain.

\section{Experimental Methods}

\subsection{Data Collection}

Measurements were conducted from January to August 2025, spanning spring, summer, and autumn seasons. The dataset comprises \num{145784} synchronized measurements, with the majority concentrated in the June-August period. Environmental conditions covered temperature from \SIrange{15}{35}{\degreeCelsius}, humidity from \SIrange{20}{45}{\percent} RH, and pressure from \SIrange{965}{1000}{\hecto\pascal}. Interferometric fringes were sampled at approximately \SI{0.06}{\hertz} while environmental parameters were recorded at \SI{1}{\hertz}.

\subsection{Data Processing}

A GPS-disciplined timebase enabled precise timestamp alignment through linear interpolation. Quality control excluded measurements during laser unlock conditions or when fringe contrast fell below \SI{50}{\percent} of maximum. The refractive index at \SI{1762}{\nano\meter} was calculated from the interference ratio $R = N_{1762}/N_{780}$, where $N_\lambda$ represents fringe counts at each wavelength.

\subsection{Statistical Analysis}

The linear model $n = \alpha_0 + \alpha_T T + \alpha_H H + \alpha_P P$ was fitted using ordinary least squares (OLS). Heteroskedasticity and autocorrelation consistent (HAC) standard errors were computed using the Newey-West estimator with \num{20} lags. Robustness was validated through \num{2000}-sample block bootstrap resampling. Error-in-variables effects were addressed via Monte Carlo simulation with \num{2000} iterations, incorporating sensor calibration uncertainties.

\section{Results}

\subsection{Refractive Index Coefficients}

Table~\ref{tab:coefficients} presents the determined coefficients with comprehensive uncertainty estimates. The model achieves $R^2 = 0.997$ with residual standard deviation $\sigma_n = \num{1.8367e-7}$. HAC standard errors account for serial correlation in the time-series data, increasing uncertainty estimates by approximately \SI{50}{\percent} compared to OLS assumptions.

\begin{table*}[ht]
\centering
\caption{Refractive index coefficients at \SI{1762}{\nano\meter} with comprehensive uncertainty estimates.}
\begin{tabular}{lcccc}
\toprule
\textbf{Coefficient} & \textbf{Value} & \textbf{OLS SE} & \textbf{HAC SE} & \textbf{95\% CI} \\
\midrule
Constant ($\alpha_0$) & $1.0000243$ & $1.1225 \times 10^{-7}$ & $1.8534 \times 10^{-7}$ & $[1.0000241, 1.0000246]$ \\
Temperature ($\alpha_T$) & $-8.8474 \times 10^{-7}$ \si{\per\degreeCelsius} & $1.7773 \times 10^{-10}$ & $2.7967 \times 10^{-10}$ & $[-8.8513, -8.8437] \times 10^{-7}$ \\
Humidity ($\alpha_H$) & $-1.3152 \times 10^{-8}$ \si{\per\percent} & $1.5564 \times 10^{-10}$ & $2.2435 \times 10^{-10}$ & $[-1.3465, -1.2859] \times 10^{-8}$ \\
Pressure ($\alpha_P$) & $+2.5949 \times 10^{-7}$ \si{\per\hecto\pascal} & $1.1356 \times 10^{-10}$ & $1.8632 \times 10^{-10}$ & $[2.5925, 2.5973] \times 10^{-7}$ \\
\bottomrule
\end{tabular}
\label{tab:coefficients}
\end{table*}

\subsection{Humidity Sensitivity Enhancement}

The measured humidity coefficient $\alpha_H = \SI{-1.3152e-8}{\per\percent}$ shows significant enhancement compared to standard models. The Ciddor model predicts \SI{-1.537e-8}{\per\percent} (difference: \SI{+15.8}{\percent}), Edlén predicts \SI{-1.526e-8}{\per\percent} (\SI{+16.6}{\percent}), and Mathar predicts \SI{-1.653e-8}{\per\percent} (\SI{+7.7}{\percent}). This enhancement originates from proximity to water absorption lines at \SI{1761.0405}{\nano\meter} and \SI{1762.4852}{\nano\meter}. Kramers-Kronig calculations predict $\alpha_H = \SI{-1.780e-8}{\per\percent}$, consistent with the observed enhancement direction and magnitude.

\subsection{Uncertainty Analysis}

Table~\ref{tab:uncertainty} presents the complete uncertainty budget. Environmental sensor calibration dominates with $\delta n \sim 2.12\times10^{-7}$, primarily from humidity measurement uncertainty. The \SI{780}{\nano\meter} reference laser contributes $\delta n \sim 10^{-9}$ at medium time scales, while the $^{138}$Ba$^+$ reference contribution is negligible ($<10^{-12}$). Statistical uncertainties from sample correlation are addressed through HAC standard errors and bootstrap validation.

\begin{table*}[ht]
\centering
\caption{Uncertainty budget for refractive index measurement at \SI{1762}{\nano\meter}.}
\begin{tabular}{lccc}
\toprule
\textbf{Source} & \textbf{Method} & \textbf{Magnitude (1$\sigma$)} & \textbf{Dominant Timescale} \\
\midrule
Environmental sensors & EIV Monte Carlo & $2.12 \times 10^{-7}$ & All \\
Sample correlation & HAC + Bootstrap & $1.69 \times 10^{-10}$ & Medium \\
780 nm reference & Analytical & $1.00 \times 10^{-9}$ & Medium \\
Spatial mismatch & Analysis & $1.84 \times 10^{-7}$ & Short \\
\bottomrule
\end{tabular}
\label{tab:uncertainty}
\end{table*}

\subsection{Statistical Robustness}

The Durbin-Watson statistic of 1.197 indicates positive autocorrelation in OLS residuals. Prais-Winsten generalized least squares correction improves this to 2.226 while reducing the residual standard deviation to $\num{1.6820e-7}$. Variance inflation factors remain below 2.0 for all predictors, indicating minimal multicollinearity. The effective sample size after accounting for autocorrelation is $N_{\text{eff}} \approx \num{39100}$.

\section{Discussion}

The \SI{12.4}{\percent} humidity sensitivity enhancement at \SI{1762}{\nano\meter} has important implications for quantum network operation. Environmental compensation algorithms that use standard atmospheric models would introduce systematic errors exceeding $10^{-6}$ in refractive index under typical humidity variations. The quantum-traceable approach eliminates this error source by providing direct SI traceability.

The dominant uncertainty from environmental sensor calibration suggests a clear path for improvement. Metrology-grade hygrometry with \SI{0.1}{\percent} RH uncertainty could reduce the total uncertainty to $\sim 10^{-9}$. This would enable refractive index measurements with accuracy comparable to the quantum frequency references themselves.

Comparison with theoretical predictions confirms the physical origin of the enhanced humidity sensitivity. The Kramers-Kronig relation links absorption features to refractive index dispersion, with water vapor lines at \SI{1761.0405}{\nano\meter} and \SI{1762.4852}{\nano\meter} creating the observed enhancement. This wavelength-specific effect underscores the need for precise measurements at quantum-relevant wavelengths rather than relying on interpolated models.

\section*{Conclusion}

In summary, we have demonstrated a direct metrological link between a primary quantum frequency standard and a key environmental parameter. By using a single trapped $^{138}$Ba$^+$ ion to provide SI traceability for a dual-wavelength interferometer, we measured the refractive index of air at 1762 nm and derived the corresponding modified Edlén coefficients with an uncertainty anchored to the quantum standard. This work effectively bridges the previously identified traceability gap, providing the first set of refractive index parameters for a quantum network wavelength that are directly validated by the qubit platform itself.

The significance of this result extends beyond the specific case of barium-ion networks. It establishes a generalizable framework for endowing quantum systems with intrinsic environmental traceability. This represents a paradigm shift from treating the environment as an uncontrolled noise source to be suppressed, toward precisely characterizing it as a measurable quantity using the system's own quantum resources. Consequently, the performance ceiling for field-deployed quantum technologies is no longer set solely by external noise models, but can be raised through self-referenced, quantum-aware compensation.

Looking forward, this co-design principle can be directly applied to other leading qubit platforms, such as neutral atoms, quantum dots, or solid-state color centers, to establish traceability for their respective operational wavelengths. Furthermore, the methodology opens a path for quantum metrology to expand its domain beyond traditional standards, contributing to fundamental atmospheric science by providing ultra-precise, wavelength-specific data for model validation. The integration of such quantum-traceable environmental sensors into network nodes will be a critical step toward realizing scalable, long-distance quantum communication and distributed quantum computing with guaranteed performance.


\section{acknowledgments}
This research was funded by the European Research Council (ERC) under the European Union’s Horizon 2020 research and innovation programme (grant number 648330), the Deutsche Forschungsgemeinschaft (DFG, grant number SCHA 973/9-1-3017959) and the Georg H. Endress Foundation. W.W. acknowledges financial support from the QUSTEC programme, funded by the European Union’s Horizon 2020 research and innovation programme under the Marie Sklodowska-Curie (grant number 847471). T.S. acknowledges financial support from the DFG via the RTG DYNCAM 2717 and the Georg H. Endress foundation.

\section{Data Availability}

The dataset and analysis code are available at [Repository URL].

\nocite{*}

\bibliography{reference}

\end{document}